\subsection{前测度与测度}

设 T 是拓扑空间，\(\mathcal{K}\) 是 T 的紧子集按包含关系排序所成的集合。对每个 \(K\in\mathcal{K}\) ，令 \(\mathcal{M}(K;C)\) 为 K 上复测度所成的集合。对 \(\mathcal{K}\) 中使 \(K\subset L\) 的每一对 (K,L)，令 \(\iota_{KL}\) 为 \(\mathcal{M}(L;C)\) 到 \(\mathcal{M}(K;C)\) 中的映射，它使 L 上的每个测度 \(\mu\) 对应于 \(\mu_{K}\) ，即 \(\mu\) 在 K 上诱导的测度 (Ch. IV, §5, No.~7, Def. 4)。则 \(\iota_{KM}=\iota_{KL}\circ\iota_{LM}\) ，只要 K、L、M 是 T 的满足 \(K\subset L\subset M\) 的紧子集；这由诱导测度的传递性得出 (Ch. V, §7, No.~2, Prop. 4)。族 \((\mathcal{M}(K;C))_{K\in\mathcal{K}}\) 关于映射 \(\iota_{KL}\) 的逆极限的元素称为 T 上的前测度。换言之：

定义 3. --- 称拓扑空间 T 上的这样的映射 w 为前测度：它使每个紧子集 K 对应 K 上的一个测度 \(w_{K}\) ，并且具有下列性质：

若 K 与 L 是 T 的使 \(K \subset L\) 的紧子集，则测度 \((w_{\mathrm{L}})_{\mathrm{K}}\) ——即 \(w_{L}\) 在 K 上诱导的测度——等于 \(w_{K}\) 。

若所有测度 \(w_{K}\) 都是实（相应地正）的，则称前测度 w 是实（相应地正）的。

设 \(w\) 与 \(w'\) 是 \(T\) 上的两个前测度，\(t\) 是一个复数；前测度 \(w + w'\) 与 \(tw\) 由公式 \((w + w')_K = w_K + w'_K\) 、\((tw)_K = tw_K\) （\(K\) 是 \(T\) 的任意紧子集）定义。\(T\) 上的前测度显然构成一个向量空间，记作 \(\mathcal{P}(T; \mathbf{C})\) ；实前测度所成的空间记作 \(\mathcal{P}(T; \mathbf{R})\) ，或更常记作 \(\mathcal{P}(T)\) ，正前测度的凸锥记作 \(\mathcal{P}_+(T)\) 。设 \(w\) 是一个前测度；则映射 \(K \mapsto |w_K|\) 是 \(T\) 上的一个前测度 (Ch. IV, §5, No.~7, Lemma 3)，记作 \(|w|\) 。若 \(w\) 是实的，则令 \(w^{+} = \frac{1}{2} (|w| + w)\) ，\(w^{-} = \frac{1}{2} (|w| - w)\) ；这两个前测度都是正的，于是可见每个实前测度都是两个正前测度之差。显然，对 T 的每个紧子集 K 有 \((w^{+})_{\mathrm{K}} = (w_{\mathrm{K}})^{+}\) ，\((w^{-})_{\mathrm{K}} = (w_{\mathrm{K}})^{-}\) 。

向量空间 \(\mathcal{P}(\mathrm{T})\) 由锥 \(\mathcal{P}_{+}(\mathrm{T})\) 赋序。显然 \(w^{+} = \sup(w,0)\) ，\(w^{-} = \sup(-w,0)\) ；因而 \(\mathcal{P}(\mathrm{T})\) 是格序的，且 \(\sup(w,w') = w + (w' - w)^{+}\) ，\(\inf(w,w') = w - (w' - w)^{-}\) 。此外，显然

\[
\big (\sup (w, w ^ {\prime}) \big) _ {\mathbf {K}} = \sup (w _ {\mathbf {K}}, w _ {\mathbf {K}} ^ {\prime}), \quad \big (\inf (w, w ^ {\prime}) \big) _ {\mathbf {K}} = \inf (w _ {\mathbf {K}}, w _ {\mathbf {K}} ^ {\prime})
\]

对 T 的每个紧子集 K 成立。

定义 4. --- 设 w 是 T 上的一个正前测度。对每个函数 \(f \in \mathcal{F}_{+}(T)\) ，我们令

\[
w ^ {\bullet} (f) = \sup _ {K} \left(w _ {K}\right) ^ {\bullet} \left(f _ {K}\right),\tag{1}
\]

其中 K 取遍 T 的紧子集所成的集合。

对每个紧集 K，令 \(p^{K}\) 为负担 \((w_{\mathrm{K}})^{\bullet}\) 在 K 到 T 中的典范单射下的像负担；\(w^{\bullet}\) 是诸负担 \(p^{K}\) 的上包络，从而是一个负担 (No.~1, Prop. 1)。称 \(w^{\bullet}\) 为与正前测度 w 相伴的本质上积分。我们常写 \(\int^{\bullet} f dw\) 或 \(\int^{\bullet} f(t) dw(t)\) 以代替 \(w^{\bullet}(f)\) 。

注 1). --- 若 \(v\) 与 \(w\) 是两个正前测度，则 \((v + w)^{\bullet} = v^{\bullet} + w^{\bullet}\) (Ch. V, §1, No.~1, Prop. 3)。若 \(v\) 与 \(w\) 是两个复前测度，则 \(|v + w|^{\bullet} \leqslant |v|^{\bullet} + |w|^{\bullet}\) 。

命题 2. --- a) 设 w 是一个正前测度。对 T 的每个紧子集 K，负担 \((w^{\bullet})_{\mathrm{K}}\) （由 \(w^{\bullet}\) 在 K 上诱导）等于 \((w_{\mathrm{K}})^{\bullet}\) 。对每个函数 \(f \in \mathcal{F}_{+}(\mathrm{T})\) ，有关系式 \((w_{\mathrm{K}})^{\bullet}(f_{\mathrm{K}}) = w^{\bullet}(f\varphi_{\mathrm{K}})\) 及

\[
w ^ {\bullet} (f) = \sup _ {K} w ^ {\bullet} (f \varphi_ {K}).\tag{2}
\]

\begin{enumerate}
\def\labelenumi{\alph{enumi})}
\setcounter{enumi}{1}
\tightlist
\item
  反之，设 \(p\) 是 \(T\) 上满足下列条件的负担：
\end{enumerate}

\begin{enumerate}
\def\labelenumi{\arabic{enumi})}
\item
  对 T 的每个紧子集 K，存在 K 上的一个正测度 \(w_{\mathbf{K}}\) 使 \(p_{\mathbf{K}} = (w_{\mathbf{K}})^{\bullet}\) 。
\item
  对每个函数 \(f \in \mathcal{F}_{+}(\mathrm{T})\) ，\(p(f) = \sup_{\mathrm{K}} p(f\varphi_{\mathrm{K}})\) 。
\end{enumerate}

则映射 \(w: \mathbf{K} \mapsto w_{\mathbf{K}}\) 是 \(\mathbf{T}\) 上的一个正前测度，且 \(p = w^{\bullet}\) 。

我们来证 \(a\) ：设 \(g \in \mathcal{F}_{+}(K)\) ，\(g^{0}\) 是 \(g\) 到 \(T\) 的 0 延拓；则由诱导负担的定义，

\[
(w ^ {\bullet}) _ {K} (g) = w ^ {\bullet} \left(g ^ {0}\right) = \sup _ {L} \left(w _ {L}\right) ^ {\bullet} \left(g ^ {0} \mid L\right),
\]

其中 L 取遍 T 的紧子集所成的集合，或仅仅取遍其中包含 K 的那些。但若 L 包含 K，则 \((w_{\mathrm{L}})^{\bullet}(g^{0}|\mathrm{L})=(w_{\mathrm{K}})^{\bullet}(g)\) ，因为 \(g^{0}|L\) 在 K 之外为零 (Ch. V, §7, No.~1, Prop. 1)，这就证明了第一个断言。于是

\[
(w _ {K}) ^ {\bullet} (f _ {K}) = (w ^ {\bullet}) _ {K} (f _ {K}) = w ^ {\bullet} \left((f _ {K}) ^ {0}\right) = w ^ {\bullet} (f \varphi_ {K})
\]

对一切 \(f \in \mathcal{F}_{+}(\mathrm{T})\) 成立，而 (2) 不过是把公式 (1) 转述了一遍。

转到 \(b\) ：1) 中所考虑的测度 \(w_{\mathbf{K}}\) 是唯一的 (Ch. V, §1, No.~1)。我们来证映射 \(\mathbf{K} \mapsto w_{\mathbf{K}}\) 是一个前测度：设 \(\mathbf{K}\) 与 \(\mathbf{L}\) 是使 \(\mathbf{K} \subset \mathbf{L}\) 的两个紧子集，\(\lambda\) 是 \(w_{\mathbf{L}}\) 在 \(\mathbf{K}\) 上诱导的测度；一切归结为证明 \(\lambda^{\bullet} = (w_{\mathbf{K}})^{\bullet}\) 。而 \(\lambda^{\bullet} = ((w_{\mathbf{L}})^{\bullet})_{\mathbf{K}}\) (Ch. V, §7, No.~1, Prop. 1)；由于 \((w_{\mathbf{L}})^{\bullet} = p_{\mathbf{L}}\) ，故 \(\lambda^{\bullet} = (p_{\mathbf{L}})_{\mathbf{K}} = p_{\mathbf{K}} = (w_{\mathbf{K}})^{\bullet}\) 。

用 \(w\) 表示前测度 \(K \mapsto w_K\) ；由于 \(p_K = (w_K)^{\bullet} = (w^{\bullet})_K\) ，有 \(p(f\varphi_K) = p_K(f_K) = (w^{\bullet})_K(f_K) = w^{\bullet}(f\varphi_K)\) 。于是两个负担 \(p\) 与 \(w^{\bullet}\) 相等，这由公式 (2) 以及关于 \(p\) 的假设 2) 得出。

证毕。

由于诱导负担 \((w^{\bullet})_{\mathbf{K}}\) 等于 \((w_{\mathbf{K}})^{\bullet}\) ，直接写 \(w_{\mathbf{K}}^{\bullet}\) 不会产生歧义。今后我们将采用这一记号。

推论. --- 设 v 与 w 是 T 上的两个正前测度，使得对 T 的每个紧子集 L 有 \(v^{\bullet}(\mathbf{L}) = w^{\bullet}(\mathbf{L})\) ；则 v = w。特别地，关系 \(v^{\bullet} = w^{\bullet}\) 蕴含 v = w。

因为，设 K 是 T 中的一个紧集；对每个紧集 \(L \subset K\) ，有关系式

\[
w _ {\mathrm{K}} (\mathrm{L}) = w _ {\mathrm{K}} ^ {\bullet} (\mathrm{L}) = w ^ {\bullet} (\mathrm{L}) = v ^ {\bullet} (\mathrm{L}) = v _ {\mathrm{K}} ^ {\bullet} (\mathrm{L}) = v _ {\mathrm{K}} (\mathrm{L})
\]

（由 Prop. 2）；故 \(w_{\mathrm{K}} = v_{\mathrm{K}}\) (Ch. IV, §4, No.~10, Prop. 19 的 Cor. 3)，最后由前测度的定义得 \(w = v\) 。

定义 5. --- 设 w 是拓扑空间 T 上的一个前测度。若负担 \(|w|^{\bullet}\) 是局部有界（相应地有界）的 (cf.~No.~1, Def. 2)，则称 w 是一个测度（相应地有界测度）。

\(\mathbf{T}\) 上复测度所成的集合显然是一个向量空间（注 1），记作 \(\mathcal{M}(\mathbf{T};\mathbf{C})\) 。实测度所成的空间记作 \(\mathcal{M}(\mathrm{T};\mathbf{R})\) ，或更常记作 \(\mathcal{M}(\mathrm{T})\) ，正测度的锥记作 \(\mathcal{M}_{+}(\mathrm{T})\) 。

若 w 是一个复测度，则其实部与虚部都是实测度。若 w 是一个实测度，则 \(w^{+}\) 与 \(w^{-}\) 都是正测度。于是每个复（相应地实）测度都是正测度的一个线性组合（相应地差）。

注记. --- 2) 若 T 局部紧，则 T 上每个前测度 w 都是测度。因为，每个 \(x \in T\) 都具有紧邻域 K，且 \(|w|^{\bullet}(K) = \|w_{K}\| < +\infty\) ，故负担 \(|w|^{\bullet}\) 局部有界。

\begin{enumerate}
\def\labelenumi{\arabic{enumi})}
\setcounter{enumi}{2}
\tightlist
\item
  对 T 的每个 Borel 子集 A（特别地对 A = T）及 T 上的每个正测度 \(\mu\) ，数 \(\mu^{\bullet}(\mathrm{A})\) 是 A 的紧子集的测度 \(\mu^{\bullet}(\mathrm{K})\) 的上确界。事实上，对 A 的每个紧子集 K 有 \(\mu^{\bullet}(\mathrm{K}) \leqslant \mu^{\bullet}(\mathrm{A})\) ；另一方面，若 \(\mathfrak{K}\) 是 T 的紧子集所成的集合，则
\end{enumerate}

\[
\mu^{\bullet}(\mathrm{A}) = \sup_{\mathrm{K}\in \mathfrak{K}}\mu_{\mathrm{K}}^{\bullet}(\mathrm{A}\cap \mathrm{K}) = \sup_{\mathrm{K}\in \mathfrak{K}}\sup_{\substack{\mathrm{L}\in \mathfrak{K}\\ \mathrm{L}\subset \mathrm{A}\cap \mathrm{K}}}\mu_{\mathrm{K}}^{\bullet}(\mathrm{L})\leqslant \sup_{\substack{\mathrm{L}\in \mathfrak{K}\\ \mathrm{L}\subset \mathrm{A}}}\mu^{\bullet}(\mathrm{L})
\]

（由 Ch. IV, §4, No.~6 的 Th. 4 的 Cor. 1）。

